% ===================== ANEXOS =====================
\appendix
\section{Cálculos Detallados}
\label{appendix:calculations}

\subsection{Cálculo del Caudal de Ventilación}

Aplicando la Ec. (\ref{eq:caudal}) con $V = 320$ m$^3$ y $N = 12$ ren/h:

\begin{equation}
	Q = \frac{320 \times 12}{60} = 64.0 \text{ m}^3/\text{min}
\end{equation}

La conversión a m$^3$/h y CFM es:

\begin{equation}
	Q = 64.0 \times 60 = 3\,840 \text{ m}^3/\text{h}
\end{equation}

\begin{equation}
	Q = 64.0 \times 35.3147 = 2\,260.1 \text{ CFM}
\end{equation}

\subsection{Cálculo de la Potencia del Ventilador}

El caudal en unidades SI es:

\begin{equation}
	Q_s = \frac{64.0}{60} = 1.0667 \text{ m}^3/\text{s}
\end{equation}

Para el escenario de diseño con $\Delta P = 250$ Pa y $\eta = 0.60$:

\begin{equation}
	P = \frac{1.0667 \times 250}{0.60 \times 1000} = 0.444 \text{ kW}
\end{equation}

En unidades de caballos de fuerza:

\begin{equation}
	P = 0.444 \times 1.341 = 0.596 \text{ HP}
\end{equation}

\subsection{Cálculo del Área de Impulsión}

Con $v_{vent} = 8$ m/s:

\begin{equation}
	A_{vent} = \frac{1.0667}{8} = 0.1333 \text{ m}^2
\end{equation}

\begin{equation}
	D = \sqrt{\frac{4 \times 0.1333}{\pi}} = 0.412 \text{ m} = 412 \text{ mm}
\end{equation}

\begin{equation}
	r = \frac{412}{2} = 206.0 \text{ mm}
\end{equation}

\subsection{Cálculo de las Rejillas de Exfiltración}

Con $v_{salida} = 3.0$ m/s:

\begin{equation}
	A_{exfil} = \frac{1.0667}{3.0} = 0.3556 \text{ m}^2
\end{equation}

Para tres rejillas:

\begin{equation}
	A_{rejilla} = \frac{0.3556}{3} = 0.1185 \text{ m}^2
\end{equation}

Seleccionando una relación de aspecto $w/h = 0.95$:

\begin{equation}
	h = \sqrt{\frac{0.1185}{0.95}} = 0.3532 \text{ m} = 353.2 \text{ mm}
\end{equation}

\begin{equation}
	w = \frac{0.1185}{0.3532} = 0.3355 \text{ m} = 335.5 \text{ mm}
\end{equation}

\subsection{Verificación del Balance de Masas}

Caudal de entrada:

\begin{equation}
	Q_{entrada} = A_{vent} \cdot v_{vent} = \frac{Q_s}{v_{vent}} \cdot v_{vent} = Q_s = 1.0667 \text{ m}^3/\text{s}
\end{equation}

Caudal de salida:

\begin{equation}
	Q_{salida} = 3 \cdot A_{rejilla} \cdot v_{salida} = 3 \times 0.1185 \times 3.0 = 1.0665 \text{ m}^3/\text{s}
\end{equation}

La discrepancia entre entrada y salida es 0.02 \%, inferior al 0.1 \%, atribuible al redondeo de dimensiones.

\section{Planos y Esquemas}
\label{appendix:drawings}

Los planos de ubicación del ventilador, rejillas de exfiltración y sensor de presión diferencial se incluirán en una entrega posterior, una vez validada la distribución de flujo mediante el modelo CFD.

\section{Especificaciones Técnicas}
\label{appendix:specs}

Las especificaciones técnicas del ventilador, filtros y rejillas se definirán en la etapa de ingeniería de detalle, con base en los parámetros de diseño entregados en el presente informe.
